% SEGMENT OLP-0599-S01
\documentclass[../../../include/open-logic-section]{subfiles}

\begin{document}

\olfileid{sth}{choice}{banach}
\olsection{பனாக்--டார்ஸ்கி முரண்புதிர்}

பூலோஸ் மாற்றீட்டை நியாயப்படுத்த முயன்றதைப் போல,
தேர்வையும் \emph{வெளிப்புற} கருத்தாய்வுகளால் நியாயப்படுத்த
முயலலாம் (\olref[replacement][extrinsic]{sec} ஐக் காண்க).
தேர்வை ஏற்றுக்கொள்வதால் பல விரும்பத்தக்க விளைவுகள்
கிடைக்கின்றன: எல்லாக் கண அளவுகளையும் ஒப்பிடலாம்;
ஒவ்வொரு கணத்தையும் நன்கு வரிசைப்படுத்தலாம்; கண அளவெண்களை
ஒருவகை வரிசையெண்களாகக் கருதலாம்; இன்னும் பல.

ஆனால் தேர்வு \emph{விரும்பத்தகாத} விளைவுகளையும் தருகிறது
என்று சில நேரங்களில் கூறப்படுகிறது. பெரும்பாலும் இதற்குக்
காரணம் \cite{BanachTarski1924} இன் முடிவு.

\begin{thm}[பனாக்--டார்ஸ்கி முரண்புதிர் ($\ZFC$ இல்)]
எந்த உருண்டையையும் முடிவுறு எண்ணிக்கையிலான துண்டுகளாகப்
பிரித்து, அவற்றைச் சுழற்றியும் இடம் மாற்றியும் மீண்டும்
அமைத்தால், அந்த உருண்டையின் இரு பிரதிகளை உருவாக்கலாம்.
\end{thm}
\noindent

% SEGMENT OLP-0599-S02
முதல் பார்வையில் இது வியப்பாக இருக்கிறது. இரு
உருண்டைகளின் கனஅளவு அசல் உருண்டையின் கனஅளவைவிட
\emph{இருமடங்கு} என்பது தெளிவு. ஆனால் உருமாறா
இயக்கங்களான சுழற்சியும் இடப்பெயர்ச்சியும் கனஅளவை
மாற்றுவதில்லை. எனவே பனாக்--டார்ஸ்கி முடிவு புதிதாகப்
பொருளை மந்திரம்போல் உண்டாக்குவதாகத் தோன்றுகிறது.

இன்னும் வியப்பானது: இதேபோன்ற காரணவாதம், ஒரு
பட்டாணியை முடிவுறு எண்ணிக்கையிலான துண்டுகளாகப் பிரித்து,
அவற்றைச் சுழற்றியும் இடம் மாற்றியும் பிக் பென் கோபுரத்தின்
வடிவமும் அளவும் கொண்ட பொருளாக அமைக்கலாம் என்பதைக்
காட்டுகிறது.\footnote{\citet[Theorem 3.12]{Wagon2016} ஐக் காண்க.}

ஆனால் இவற்றில் எதையும் $\ZF$ மட்டும் கொண்டு பெற
முடியாது.\footnote{தேர்வைவிடக் கண்டிப்பாக வலுக்குறைந்த
கொள்கைகளாலும் பனாக்--டார்ஸ்கி முடிவை நிறுவலாம்;
\citet[303]{Wagon2016} ஐக் காண்க.} ஆகவே நாம் ஒரு
முடிவெடுக்க வேண்டும்: தேர்வை மறுப்பதா, அல்லது இந்த
``முரண்புதிரை'' ஏற்றுக்கொள்வதா?

% SEGMENT OLP-0599-S03
இந்த ``முரண்புதிருடன்'' வாழக் கற்றுக்கொள்ளலாம் என்பதே
எங்களது கருத்து. உண்மையில், இது பெரிய முரண்புதிராகவே
எங்களுக்குத் தோன்றவில்லை. குறிப்பாகப் பின்வரும்
முடிவுகளைவிட இது ஏன் அதிகமாகவோ குறைவாகவோ
முரண்பாடாகத் தோன்ற வேண்டும்?\footnote{வேறு எடுத்துக்காட்டுகளால்
\citet[276--7]{Potter2004}, \citet[16]{Weston2003},
\citet[31, 308--9]{Wagon2016} ஆகியோரும் இதைப் போன்ற
கருத்துகளை முன்வைக்கின்றனர்.}
\begin{enumerate}
	\item $(0,1)$ இடைவெளியிலுள்ள புள்ளிகளின் எண்ணிக்கையும்
	$\Real$ இல் உள்ள புள்ளிகளின் எண்ணிக்கையும் ஒன்றே.
	% SOURCE CORRECTION TA-STH-046: டேன்ஜென்ட் வாய்பாட்டின் கூடுதல் மூடும் அடைப்பை நீக்குக.
	\\\emph{நிறுவல்}: $\tan(\pi(r-\nicefrac{1}{2}))$ ஐக் கருதுக.
	\item ஒரு கோட்டிலுள்ள புள்ளிகளின் எண்ணிக்கையும் ஒரு
	சதுரத்திலுள்ள புள்ளிகளின் எண்ணிக்கையும் ஒன்றே.
	\\\olref[his][set][pathology]{sec} மற்றும்
	\olref[his][set][cantorplane]{sec} ஐக் காண்க.
	\item இடத்தை நிரப்பும் வளைவுகள் உள்ளன.
	\\\olref[his][set][pathology]{sec} மற்றும்
	\olref[his][set][hilbertcurve]{sec} ஐக் காண்க.
\end{enumerate}
இந்த மூன்று முடிவுகளுக்கும் தேர்வு தேவையில்லை. இப்போது
இவற்றை வியப்பூட்டும், அழகிய கணிதமாகவே கருதுகிறோம்.
பனாக்--டார்ஸ்கி முரண்புதிரையும் அப்படியே கருதலாம்.

% SEGMENT OLP-0599-S04
இங்கு ஒரு நுட்பமான தொழில்நுட்பக் குறிப்புத் தேவை;
ஆனால் நிதானமாக நினைத்தால் போதும். உருமாறா இயக்கங்கள்
கனஅளவைப் பாதுகாக்கின்றன. ஆகவே உருண்டை பிரிக்கப்படும்
ஐந்து\footnote{முரண்புதிரை ``முடிவுறு எண்ணிக்கையிலான
துண்டுகள்'' என்று கூறினோம். உண்மையில், \citet{Robinson1947}
\emph{ஐந்து} துண்டுகள் போதும், அதைவிடக் குறைய முடியாது
என்று நிறுவினார். ஒரு நிறுவலுக்கு \citet[pp.~66--7]{Wagon2016}
ஐப் பார்க்கவும்.} துண்டுகளும் \emph{அளவிடத்தக்கவையாக}
இருக்க முடியாது. எளிமையாகச் சொன்னால், ஒவ்வொரு
துண்டுக்கும் தனித்தனியாக ஒரு கனஅளவை ஒதுக்குவது
பொருளற்றது. அவற்றைப் படமாகக் காண முடியாத,
``முடிவின்றிச் சிதறிய'' புள்ளிகளின் தொகுப்புகளாகக் கருதலாம்.
இவ்வாறு ``முடிவின்றிச் சிதறிய'' கணங்கள் இருப்பது
விசித்திரமாகத் தோன்றலாம். ஆனால் முன்பே கிடைக்கும்
கணங்களின் \emph{எல்லாச் சாத்தியமான} தொகுப்புகளையும்
உருவாக்க வேண்டும் என்ற \stagesacc{} கொள்கையிலிருந்து
அவை கிடைப்பதாகத் தோன்றுகின்றன.

% SEGMENT OLP-0599-S05
இவை எதுவும் உங்களை நம்பவைக்கவில்லை எனில்,
பனாக்--டார்ஸ்கி முரண்புதிரை ஏற்றுக்கொள்வதற்கான
இறுதி (வெளிப்புற) காரணம் ஒன்று இருக்கிறது. அது
சிறந்த கணித நகைச்சுவையை உடனே தருகிறது:
\begin{enumerate}
	\item[] \emph{கேள்வி}. ``பனாக்--டார்ஸ்கி'' என்ற
	பெயரின் எழுத்துகளை மறுவரிசைப்படுத்தினால் என்ன கிடைக்கும்?
	\item[] \emph{பதில்}. ``பனாக்--டார்ஸ்கி பனாக்--டார்ஸ்கி''.
\end{enumerate}

\end{document}
